\section{Exact operator bridge to the released Navier--Stokes profiles}
\label{nsb:operator}

This bounded comparison uses the locally saved released Navier--Stokes source
(the source pages are 7--15 and 24--27 for the profile, stress and pulse
construction). Its stated physical viscosity is denoted by $\nu_{\rm NS}>0$.
The Boussinesq lane variables and physical viscosity remain $d$, $\mu$ and
$\nu=\sqrt{A_0}$; none is identified with $\nu_{\rm NS}$.

Use right-handed cylindrical coordinates $(r,\vartheta,z)$ and define
\[
 y_1=z,\qquad y_2=\frac{r^2}{2}>0,\qquad
 v=(u^z,ru^r)=J\nabla_y\Psi,
 \quad J=\begin{pmatrix}0&-1\\1&0\end{pmatrix}.
\]
The meridional divergence identity is
$\partial_{y_1}v_1+\partial_{y_2}v_2=0$. Define the circulation and
meridional vorticity variables
\[
 \Gamma=ru^\vartheta,
 \qquad \xi=\frac{\omega^\vartheta}{r},
 \qquad h=\frac{\Gamma}{2y_2}=\frac{u^\vartheta}{r},
 \qquad Q=\partial_{y_1}^2+2y_2\partial_{y_2}^2,
 \qquad M=Q+4\partial_{y_2}.
\]
The azimuthal momentum and meridional vorticity rows of the cylindrical
equations (with arbitrary smooth force components) become, on $y_2>0$,
\begin{align}
 D_v\Gamma&=\nu_{\rm NS}Q\Gamma+r f^\vartheta,
 \label{nsb:Gamma}\\
 D_v\xi&=\nu_{\rm NS}M\xi+\partial_{y_1}(h^2)
       +\operatorname{curl}_y(f^z,f^r/r),
 \label{nsb:xi}
\end{align}
where $D_v=\partial_t+v\cdot\nabla_y$ and
$\operatorname{curl}_y(a,b)=\partial_{y_1}b-\partial_{y_2}a$.

Indeed, $\partial_r=r\partial_{y_2}$ and
$\partial_r^2+r^{-1}\partial_r=2y_2\partial_{y_2}^2+2\partial_{y_2}$.
For $\Gamma$, the azimuthal viscous operator is
$\partial_r^2-r^{-1}\partial_r+\partial_z^2=Q$.
For $r\xi=\omega^\vartheta$, the cylindrical operator
$(\partial_r^2+r^{-1}\partial_r-r^{-2}+\partial_z^2)(r\xi)$,
divided by $r$, is $M\xi$; explicitly the radial terms are
$2y_2\xi_{y_2y_2}+4\xi_{y_2}$. The centrifugal term is
$r^{-2}\partial_z((u^\vartheta)^2)=\partial_{y_1}(\Gamma^2/(4y_2^2))
=\partial_{y_1}(h^2)$. These calculations prove every coefficient in
(\ref{nsb:Gamma})--(\ref{nsb:xi}), including the $r$ force conversion.

The velocity--vorticity relation is also retained. Since
$u^z=-\Psi_{y_2}$ and $u^r=\Psi_{y_1}/r$,
\[
 \xi=\frac{\partial_z u^r-\partial_r u^z}{r}
     =\frac{\Psi_{y_1y_1}}{2y_2}+\Psi_{y_2y_2}
     =\frac{Q\Psi}{2y_2}.
\]
Thus the two transport rows do not by themselves replace the elliptic
velocity reconstruction. Conversely, on a simply connected meridional
domain, a smooth $\Psi,\Gamma$ satisfying this relation and both transport
rows recovers $u$. Define its momentum residual without pressure as
$a=\partial_tu+(u\cdot\nabla)u-\nu_{\rm NS}\Delta u-f$.
The circulation row is $a^\vartheta=0$, while the vorticity row is
$\partial_z a^r-\partial_r a^z=0$. The one-form
$a^r\,\mathrm dr+a^z\,\mathrm dz$ is therefore closed, and
$p=-\int(a^r\,\mathrm dr+a^z\,\mathrm dz)$ along a path from a fixed
base point is path independent. Its derivatives are $p_r=-a^r$,
$p_z=-a^z$. This proves the local momentum reconstruction, with pressure
unique up to a function of time; boundary data are not removed.

The change from circulation to swirl rate is invertible on the exact retained
domain: $\Gamma=2y_2h$. Product differentiation gives
\[
 Q(2y_2h)=2y_2Mh,\qquad
 D_v(2y_2h)=2y_2D_vh+2v_2h.
\]
Therefore (\ref{nsb:Gamma}) is equivalent to
\begin{equation}
 D_vh=\nu_{\rm NS}Mh-\frac{v_2}{y_2}h+\frac{f^\vartheta}{r}.
 \label{nsb:h}
\end{equation}
The source profile uses precisely this relation: in its similarity variables,
$E$ is the azimuthal velocity profile and
$H_{\rm src}=\sqrt{2X}E$ is the scaled circulation profile (source equations
(4.3), (4.8)). Since $r=\sqrt{2qX}$ and
$u^\vartheta=q^{-A}E$, its physical circulation is
$\Gamma=q^{1/2-A}H_{\rm src}=q^{-h_{\rm src}}H_{\rm src}$,
where the source exponents satisfy $A=1/2+h_{\rm src}$. Here $q>0$ is the
geometric similarity coordinate and $h_{\rm src}>0$ is an exponent; neither
is a signed amplitude scale. This is the exact scale dictionary; it does not
identify $H_{\rm src}$ with the Boussinesq scalar amplitude.

The sign of the blowout amplitude must be kept separate from that positive
geometric coordinate. The axisymmetric equations have the exact involution
\[
 (u^r,u^\vartheta,u^z,p;f^r,f^\vartheta,f^z)
 \longmapsto
 (u^r,-u^\vartheta,u^z,p;f^r,-f^\vartheta,f^z).
 \label{nsb:sign}
\]
It sends $(\Gamma,h)$ to $(-\Gamma,-h)$, leaves $\xi$ and the meridional
flow unchanged, and preserves every displayed equation: the only quadratic
swirl term is $\partial_{y_1}(h^2)$, while the azimuthal equation is linear in
$(\Gamma,f^\vartheta)$. Consequently a source-selected positive profile
$E>0$ (source pp. 7--9) does not justify assuming that every blowout branch is
positive; its reflected branch has $E<0$, $\Gamma<0$ and the same blowup
magnitude. The exponent $q^{-h_{\rm src}}$ controls the magnitude in either
branch. The source's positive covariance coefficients and viscosity are
separate construction choices, not a sign claim about the velocity amplitude.

For comparison with the lane's controlled scalar equation, set
$B=h^2$. The globally smooth square equation, including its geometric and
viscous terms, is
\begin{equation}
 D_vB=\nu_{\rm NS}MB-2\nu_{\rm NS}
 |\nabla h|_Q^2-\frac{2v_2}{y_2}B+\frac{2h}{r}f^\vartheta,
 \qquad |\nabla h|_Q^2=h_{y_1}^2+2y_2h_{y_2}^2.
 \label{nsb:B}
\end{equation}
This follows by multiplying (\ref{nsb:h}) by $2h$ and using
$M(h^2)=2hMh+2|\nabla h|_Q^2$. At zeros of $h$ equation
(\ref{nsb:B}) remains smooth; dividing by $B$ is unnecessary and would lose
that domain. On a component where $B>0$, the equivalent expression is
$D_vB=\nu_{\rm NS}MB-\nu_{\rm NS}|\nabla B|_Q^2/(2B)
-2v_2B/y_2+2hf^\vartheta/r$.

The exact relation to the Boussinesq lane is consequently an operator
correspondence, not a scalar conjugacy. The lane's retained-viscosity scalar
row has $D_v\Theta=d\Delta\Theta+f_\Theta$ in Cartesian material
coordinates. A direct identification $\Theta=B$ would require, on the same
physical domain, both $M=\Delta$ after the nonlinear coordinate pullback and
$-2v_2B/y_2-2\nu_{\rm NS}|\nabla h|_Q^2+2hf^\vartheta/r$ to be absorbed into
one prescribed force. The first condition fails as an operator identity:
$M=\partial_{y_1}^2+2y_2\partial_{y_2}^2+4\partial_{y_2}$ has variable metric
and drift, whereas the lane's $\Delta$ is the Euclidean Cartesian Laplacian.
The geometric, gradient and force terms must therefore be retained or
calculated in the particular source state. Nonconstant profiles alone
would not prove that such terms never cancel.
The exact bridge is (\ref{nsb:Gamma})--(\ref{nsb:B}) with the
full residual terms retained. It neither claims that the lane's controls are
the source pulses nor changes the source's pressure reconstruction.

For the meridional velocity, the source's pulse construction supplies the
following further exact local dictionary (source pp. 11--15 and 74--77). If
$w$ is a divergence-free pulse and $\pi$ its pressure increment, then
\[
 R(u_B+w,p_B+\pi)=R(u_B,p_B)+L_{u_B}(w,\pi)
 +\nabla\!\cdot(w\otimes w),
\]
\[
 L_{u_B}(w,\pi)=\partial_tw+(u_B\cdot\nabla)w
 +(w\cdot\nabla)u_B-\nu_{\rm NS}\Delta w+\nabla\pi,
 \qquad \nabla\cdot w=0.
\]
For the moving local phase with wavevector $n$ and transverse amplitude
$t_m$, the source equation is
\[
 t_m'+\mathcal Kt_m+m^2d_N t_m+i k m n\,\pi_m=-f_m,
 \qquad n\cdot t_m=0,
 \qquad d_N=\varepsilon k^2|n|^2,
\]
with the source's moving shear matrix
$\mathcal K=\begin{pmatrix}0&-2F&0\\2F+RF_R&0&0\\G_R&0&0\end{pmatrix}$.
The transverse projection uses
\[
 A_\Phi=-\mathcal K+
 \frac{n(n^T\mathcal K-n'^T)}{|n|^2}.
\]
These are exact source pulse equations, including viscosity and pressure
projection. Their quadratic covariance is chosen to cancel the annular
stress; the source explicitly retains curl, cutoff, and correction interactions.
The Boussinesq lane's two-dimensional pair and its full profile residuals have
an exact algebraic form, but no identification with this three-dimensional
pulse equation follows without an additional embedding that supplies the
source's cylindrical geometry, moving plane, covariance cone and correction
cycle. That embedding is not asserted here.

The current bridge therefore proves a precise morphism of variables, operators,
forces and pulse residuals and records the exact obstruction to a direct scalar
identification. It establishes the next programme datum while preserving the
unproved status of any singular limit or third return.
